\centerline{\bf \Huge Квадратичные вычеты}


В этом файле мы постараемся разобраться с такой темой, как квадратичные вычеты. На самом-то деле более полную информацию об этой теме можно прочесть в прекраснейшей \href{https://drive.google.com/file/d/1eIYNdSgwb4Nw8PnFT2Ng2LWQURzPk129/view?usp=sharing}{книжке} Павла Викторовича Бибикова. Я лишь постараюсь вкратце изложить суть, которой будет достаточно, чтобы начать решать задачи.

Перед тем как начать, вспомним некоторые важные факты из теории чисел:

\begin{enumerate}
    \item \textit{Малая теорема Ферма}. $a^{p-1} \equiv 1 \pmod{p}$ при $a \not\equiv 0 \pmod{p}$
    \item \textit{Теорема Эйлера}. $a^{\varphi(n)} \equiv 1 \pmod{n}$ при $(a,n) = 1$
    \item теорема Вильсона: $(p-1)!\equiv -1 \pmod{p}$
    \item Через $\mathbb{Z}_n^*$ обозначается множество всех остатков под модулю $n$ взаимно простых с $n$. Так же это называется мультипликативной группой вычетов по модулю $n$ или же приведенной системой вычетов. В частности $\mathbb{Z}_p^* = \{1, 2, 3, \ldots, p-1\}$, оно также называется полем вычетом по модулю $p$.
    \item Сравнение $ax = 1 \pmod{n}$ имеет решение (в переменной $x$) тогда и только тогда, когда $(a,n) = 1$.
\end{enumerate}

\textbf{Определение 1.}
Остаток а называется \textit{квадратичным вычетом} (или просто вычетом) по модулю $p$, если существует такое целое число $x$, что $x^2 \equiv_p a$. В противном случае число $a$ называется \textit{квадратичным невычетом} (или просто \textit{невычетом}).


\textbf{Утверждение.}
Для нечетного простого числа $р$ количество квадратичных вычетов в поле $\mathbb{Z}_p$ равно $\frac{p-1}{2}$. Столько же будет и невычетов.

\textit{Доказательство.} Заметим, что остатки $x$ и $p-x$, не равные 0 , при возведении в квадрат дают одинаковые остатки: $x^2 \equiv(p-x)^2 \pmod{p}$. Поэтому квадратичных вычетов не более $\frac{p-1}{2}$ : мы отбросили остаток 0 , а все остальные остатки разбили на пары ( $x, p-x$ ), квадраты которых дают одинаковые остатки, являющиеся квадратичными вычетами. Докажем, что все эти вычеты будут различны. В самом деле, если $x^2 \equiv y^2 \pmod{p}$, то, перенося $y^2$ в левую часть, получаем, что $(x-y)(x+y) \equiv 0 \pmod{p}$. Это означает, что произведение $(x-y)(x+y)$ делится на простое число $p$. Это возможно в одном из двух случаев: либо $x-y$ делится на $p$ (и тогда $y \equiv x\pmod{p}$ ), либо $x+y$ делится на $p$ (и тогда $y \equiv-x \pmod{p}$ ). Значит, не более двух различных остатков при возведении в квадрат могут совпасть. Но такие остатки мы знаем - это в точности пары остатков $(x, p-x)$. Каждая такая пара дает свой квадратичный вычет, и самих вычетов получается в точности $\dfrac{p-1}{2}$ что и требовалось доказать.


\textbf{Определение 2.}
Чтобы можно было обозначать свойство числа быть вычетом или невычетом, используется специальный символ, известный также как символ Лежандра. Он обозначается через $\left(\frac{a}{p}\right)$ и определяется следующим образом:

$$
\left(\frac{a}{p}\right)=\left\{\begin{array}{l}
0, \quad \text { если } a \equiv_p 0, \\
1, \quad \text { если } a-\text { вычет в } \mathbb{Z}_p, \\
-1, \quad \text { если } a-\text { невычет в } \mathbb{Z}_p .
\end{array}\right.
$$

\textbf{Теорема 1.} \textit{(Критерий Эйлера)}
Для любого простого числа $р$ и целого числа а имеет место сравнение $\left(\frac{a}{p}\right) \equiv_p a^{\frac{p-1}{2}}$.

Доказательство. Для доказательства критерия Эйлера нам достаточно проверить два утверждения:

--- если $\left(\frac{a}{p}\right)=1$, то $a^{\frac{p-1}{2}} \equiv_ 1\pmod{p}$;

--- если $\left(\frac{a}{p}\right)=-1$, то $a^{\frac{p-1}{2}} \equiv_-1\pmod{p}$.

Доказательство первого утверждения не представляет особой сложности. Действительно, если $a$ - квадратичный вычет, то найдется такой ненулевой остаток $x$, что $a \equiv_p x^2$. Тогда

$$
a^{\frac{p-1}{2}} \equiv_p\left(x^2\right)^{\frac{p-1}{2}}=x^{p-1} \equiv_ 1\pmod{p}
$$


по малой теореме Ферма.

Теперь докажем второе утверждение. Пусть $a$ --- квадратичный невычет. Рассмотрим многочлен $f(t)=t^{\frac{p-1}{2}}-1$ над полем $\mathbb{Z}_p$. По теореме Безу у этого многочлена есть не более $\frac{p-1}{2}$ корней. Но по первому утверждению все квадратичные вычеты являются корнями многочлена $f$, а вычетов, как мы помним, ровно $\frac{p-1}{2}$, т.е. их количество равно степени многочлена $f$. Значит, если $a$ --- невычет, то он не может быть корнем многочлена $f$, т.е. $f(a) \not\equiv 0 \pmod{p}$ и $a^{\frac{p-1}{2}}-1 \neq 0 \pmod{p}$. Но мы уже отмечали выше, что $a^{\frac{p-1}{2}} \equiv 1 \pmod{p}$ или -1 . Раз случай 1 невозможен, остается только вариант $с-1$. Поэтому $a^{\frac{p-1}{2}} \equiv-1 \pmod{p}$, что и требовалось доказать.


\textbf{Теорема 2.}
Для любых ненулевых остатков а и $b$ по нечетному простому модулю $p$ имеет место равенство $\left(\frac{a b}{p}\right)=\left(\frac{a}{p}\right)\left(\frac{b}{p}\right)$. Иначе говоря, символ Лежандра мультипликативен.


\textit{Доказательство.}
$$\left(\frac{a b}{p}\right) \equiv_p(a b)^{\frac{p-1}{2}}=a^{\frac{p-1}{2}} \cdot b^{\frac{p-1}{2}} \equiv_p\left(\frac{a}{p}\right)\left(\frac{b}{p}\right)$$.

\textit{Замечаение.} Отметим, что отсюда следует такой нетривиальный с виду факт:

--- произведение двух вычетов - вычет;

--- произведение вычета и невычета - невычет;

--- произведение двух невычетов - вычет.

На алгебраическом языке это означает, что множество квадратичный вычетов является подгруппой индекса 2 группы вычетов $\mathbb{Z}_p^*$.

Мы с вами прекрасно знаем критерий того есть ли решение у уравнения $ax^2 + bx + c = 0$ и если они есть, то сколько их. Если $D>0$, то корней два, если $D=0$, то корень один, а если $D<0$, то корней нет. Оказывается примерно такой же критерий есть и для решения сравнения $ax^2 + bx + c \equiv 0 \pmod{p}$.

\textbf{Теорема 3.}
Пусть $a, b$ и $c$ - произвольные целые числа, такие, что $a$ не делится на нечетное простое число $p$. Докажите, что сравнение $a t^2+bt +c \equiv 0 \pmod{p}$ имеет два решения если $\left(\dfrac{D}{p}\right) = 1$, одно решение, если $\left(\dfrac{D}{p}\right) = 0$ и не имеет решений, если $\left(\dfrac{D}{p}\right) = 0$, где $D=b^2-4 a c-$ дискриминант квадратного уравнения.

\textit{Доказательство.}
Чтобы решить сравнение, выделим полный квадрат. Для этого нам будет удобно домножить сравнение на $4 a$ :

$$
4 a^2 t^2+4 a b t+b^2 \equiv_p b^2-4 a c \quad \Longrightarrow \quad(2 a t+b)^2 \equiv_p D .
$$
Отсюда и следует требуемое. 

\textbf{Теорема 4.}
\textit{(о корне из $-1$)}
Число  $-1$ является квадратичным вычетом по нечетному простому модулю $p$ тогда и только тогда, когда $p \equiv 1\pmod{p}$, т.е. $p$ имеет вид $4 k+1$.

\textit{Доказательство.} 
По критерию Эйлера имеем $\left(\frac{-}{p}\right)=(-1)^{\frac{p-1}{2}}$. Правая часть этого равенства равна 1 тогда и только тогда, когда число $\frac{p-1}{2}$ четно, т.е. когда $p \equiv_4 1$, что и требовалось доказать.

Покажем одно из применений теоремы о корне из $-1$.

\textbf{Теорема 5.}
\textit{(теорема Жирара)}. Если целье числа $x$ и у таковы, что $x^2+y^2 \equiv_p 0$ для некоторого простого числа $p$ вида $4 k+3$, то $x \equiv_p y \equiv_p 0$.

\textit{Доказательство.} Предположим, что существуют такие $x$ и $y$, что $y \not \equiv 0 \pmod{p}$ и $x^2+y^2 \equiv 0 \pmod{p}$. Разделив это сравнение на $y^2$ и полагая $t=x y^{-1}$, запишем сравнение в виде $t^2 \equiv_p-1$. Но это сравнение невозможно в силу теоремы о корне из $-1$.

Теперь перейдём к ключевой теореме во всей этой теории. 

\textbf{Теорема 6.}
\textit{(Квадратичный закон взаимности Гаусса).}
Пусть $p$ и $q$ - различные нечётные простые числа. Тогда имеет место равенство

$$
\left(\frac{p}{q}\right) \cdot\left(\frac{q}{p}\right)=(-1)^{\frac{p-1}{2} \cdot \frac{q-1}{2}}
$$

\textit{Доказательство.} Остается в качестве упражнения читателю)))

Как вы могли заметить везде простые числа, которые фигурировали в наших утверждениях, были нечётными. А что же делать с двойкой? Оказывается верен следующий полезный факт:

\textbf{Утверждение.}
Для нечётного простого $р$ что число $2$ является квадратичным вычетом тогда и только тогда, когда $p$ имеет вид $8 k \pm 1$. Иными словами, $\left(\frac{2}{p}\right)=(-1)^{\frac{p^2-1}{8}}$. Иначе говоря 2 вычет, если $p \equiv \pm1 \pmod{8}$ и невычет, если $p \equiv \pm 3 \pmod{8}$


Теперь мы набрались полным чемоданом фактов, чтобы наконец решать задачи. Покажем их применение в действии на нескольких примерах.


\newpage

\hspace{3mm}

\centerline{\bf \Huge Примеры решения задач.}

\textbf{Пример 1.} Как же всё-таки вычислять эти символы Лежандра??? Сделаем пошаговую инструкцию, как вычислить $\left(\dfrac{a}{p}\right)$

\begin{enumerate}
    \item Если $a>p$, то заменим его на его остаток при делении на $p$.
    \item Если после первого пункта (или же до него) у вас получился квадрат натурального числа, то сразу же пишем ответ 1.
    \item Разложите $a$ на простые множители $a = p_1^{\alpha_1}\cdotp p_2^{\alpha_2} \cdot \ldots \cdot p_k^{\alpha_k}$
    \item По мультипликативности символ Лежандра перепишется вот так

    $$
    \left(\dfrac{a}{p}\right) = \left(\dfrac{p_1^{\alpha_1}\cdotp p_2^{\alpha_2} \ldots p_k^{\alpha_k}}{p}\right) = \left(\dfrac{p_1}{p}\right)^{\alpha_1}\cdot\left(\dfrac{p_2}{p}\right)^{\alpha_2}\cdot \ldots \cdot \left(\dfrac{p_k}{p}\right)^{\alpha_k}
    $$
    \item После предыдущего пункта нам лишь осталось научиться вычислять символ Лежандра $\left(\dfrac{p}{q}\right)$, где $p,q$ простые числа. И вот тут нам поможет квадратичный закон взаимности Гаусса. Если хорошо задуматься, то он говорит следующее
    $$
    \left(\frac{p}{q}\right)=\left\{\begin{array}{l}
    0, \quad \text { если } p =q, \\
    -\left(\dfrac{q}{p}\right), \quad \text { если }  p\equiv q \equiv -1\pmod{4}\\
    \left(\dfrac{q}{p}\right), \quad \text { во всех остальных случаях.} 
    \end{array}\right.
    $$

    То есть если одно из чисел $p,q$ даёт остаток 1 при делении на 4, то переворачиваем символ Лежандра без изменения знака. А если оба числа дают остаток 3, то переворачиваем с изменением знака. Так же не забываем, как вычислять символы Лежандра с числителем -1 и 2.
\end{enumerate}

Но что даёт мне это переворачивание? Давайте посмотрим на примере. Попробуем вычислить символ Лежандра $\left(\dfrac{23}{47}\right)$. Оба числа простых числа дают остаток 3 при делении на 4, поэтому $\legendre{23}{47} = -\legendre{47}{23}$. Теперь вспоминаем первый пункт и заменяем 47 на его остаток. Финальная цепочка равенств будет выглядеть вот так:

$$
\legendre{23}{47} = -\legendre{47}{23} = \legendre{1}{23} = 1
$$

Теперь приведём сложный пример:

$$
\legendre{1001}{277} = \legendre{170}{277} = \legendre{2}{277} \cdot \legendre{5}{277} \legendre{17}{277} = (-1)\cdot \legendre{277}{5} \cdot \legendre{277}{17} =
$$
$$
=(-1)\cdot \legendre{2}{5}\legendre{5}{17} = (-1)\cdot(-1)\cdot\legendre{17}{5} = (-1)(-1)\legendre{2}{5} = (-1)(-1)(-1) = -1
$$

Попробуйте самостоятельно осознать каждый переход.

Если вы думаете, что квадратичный закон взаимности помогает только лишь вычислять конкретные значения символов Лежандра, то вы ошибаетесь. Следующая задача тому подвтерждение.

\textbf{Пример 2.} 
Докажите, что простых чисел вида $10 k-1$ бесконечно много.

\textit{Доказательство.}
Докажем, что простых чисел сравнимых с $- 1$ по модулю $5$ бесконечно много. Это будет равносильно задаче в силу того, что все простые кроме числа $2$ - нечетные. Пусть это не так, то есть таких чисел конечно, обозначим за $A$ удвоенное их произведение. Число $A^2-5$ - больше 1 и нечетно, тогда $A^2-5 \equiv 0(\bmod p)$. Тогда 5 - квадратичный вычет по модулю $p$, и по квадратичному закону взаимности имеем:

$$
\left(\frac{5}{p}\right) \cdot\left(\frac{p}{5}\right)=(-1)^{\frac{p-1}{2} \cdot \frac{5-1}{2}}=1 \Rightarrow\left(\frac{p}{5}\right)=1
$$


то есть $p \equiv \pm 1 \pmod 5$. Но все простые делители не могут быть вида $5 k+1$ ведь $A^2-5 \equiv 4 \pmod 5$. Значит существует еще одно простое число вида $5 k-1$ - противоречие.

\textbf{Пример 3}
Найдите все целые числа $x$ и $y$, для которых число $\frac{x^2+x+2}{y^6-2}$ целое.

\textit{Доказательство.}
При $y= \pm 1,0$ очевидно, что $x$ может быть любым. Во всех остальных случаях $y^6-2>0$. По малой теореме Ферма $y^6-2 \equiv-1 \pmod{7}$ или -2 . Оба этих числа являются невычетами по модулю 7 , поэтому у $y^6-2$ найдется простой делитель $p$, являющийся невычетом по модулю 7 (в частности, $p \neq 2$ ). Заметим, что $(2 x+1)^2 \equiv-7 \pmod{p}$, то есть $\left(\frac{-7}{p}\right)=1$. По КЗВ имеем

$$
\left(\frac{-7}{p}\right)=\left(\frac{-1}{p}\right)\left(\frac{7}{p}\right)=(-1)^{(p-1) / 2} \cdot\left(\frac{7}{p}\right) \cdot(-1)^{(p-1)(7-1) / 2}=-1
$$


Таким образом, $x^2+x+2$ не делится на $p$, и дробь целых значений не принимает.

Ответ:
$y= \pm 1,0$ и $x$ - любое целое число

\textbf{Пример 4.}
Дано натуральное число $n>1$. Докажите, что любой нечётный делитель $d>10$ числа $5 n^2+1$ имеет в своей десятичной записи чётную цифру.

\textit{Доказательство.} Докажем, что вторая цифра числа $d$ будет четной. Для этого достаточно доказать, что $d \equiv_{20} 1,3,5,7,9$. Заметим, что случай $d \equiv_{20} 5$ невозможен, т.е. нужно проверить, что остатки делителя $d$ по модулю 20 могут быть равны лишь $1,3,7$ или 9 . Далее, заметим, что если числа $d_1$ и $d_2$ имеют такие остатки, то остаток произведения $d_1 d_2$ также будет равен 1,3 , 7 или 9: это следует из таблицы умножения по модулю 20 :

\begin{tabular}{|c|l|l|l|l|}
\hline$\times$ & 1 & 3 & 7 & 9 \\
\hline 1 & 1 & 3 & 7 & 9 \\
\hline 3 & 3 & 9 & 1 & 7 \\
\hline 7 & 7 & 1 & 9 & 3 \\
\hline 9 & 9 & 7 & 3 & 1 \\
\hline
\end{tabular}

Таким образом, достаточно доказать, что любой простой делитель $p>10$ числа $5 n^2+1$ имеет остаток $1,3,7$ или 9 по модулю 20. Для этого вычислим остатки числа $p$ по модулям 4 и 5 . Если $5 n^2 \equiv_p-1$, то $(5 n)^2 \equiv_p-5$, т.е. -5 - вычет по модулю $p$. Применим квадратичный закон взаимности:

$$
\left(\frac{-5}{p}\right)=\left(\frac{-1}{p}\right)\left(\frac{5}{p}\right)=(-1)^{\frac{p-1}{2}} \cdot\left(\frac{p}{5}\right) \cdot(-1)^{\frac{(p-1)(5-1)}{4}}=(-1)^{\frac{p-1}{2}} \cdot\left(\frac{p}{5}\right) .
$$


Значит, либо $p \equiv_4 1$ и $p \equiv_5 \pm 1$, либо $p \equiv_4-1$ и $p \equiv_5 \pm 2$. Несложно проверить, что все четыре варианта дают в точности остатки $1,3,7$ и 9 по модулю 20 . Таким образом, наша задача решена.